\section{Regular patches and the size of the conformal forcing}
\label{sec:rpi:regular-patches}

Our next objective is to separate two distinct flow frontiers, including
their singular points. At a hypothetical contact, rescaling makes both
frontiers approach the same minimizing cone. Its regular patches are
the places where their separation can be described by a scalar graph
equation. Before using that equation, we need estimates for the whole
intervening time family, and we must measure the harmonic forcing
relative to its own size. A strict Hopf derivative at each fixed scale
would not by itself survive normalization.

Fix $T<\infty$, and write $k=n-1$ and $d=n-2$. We use the flow
already constructed: $K_t$ is its closed filled side,
$E_t=M\setminus K_t$ its connected end exterior, and
$\Sigma_t=\supp|D\chi_{K_t}|$ its full perimeter support. The
normal always points from the filled side into this exterior. In the
smooth reference metric $g_0$, the frontiers have uniform local
perimeter and two-phase density bounds and an almost-minimizing error
$\beta r^{k+\mu}$ with $\mu>0$. We also have
\[
 g_t=u_t^{4/d}g_0,\qquad
 0<c_T\le u_t\le1,\qquad |v_t|\le1,\qquad
 u_t=1+\int_0^t v_\sigma\,d\sigma.
\]
The factors have a common background H\"older modulus. The sets are
nested, each $v_t$ is harmonic in $E_t$ and zero on $K_t$, and each
$u_t$ is harmonic in $E_t$. The initial frontier is locally perimeter
minimizing in $g_0$. All neighborhoods below avoid the buried physical
boundary.

The rescaled almost-minimizing errors tend to zero. Compactness,
perimeter convergence, and small-excess regularity therefore give
multiplicity-one minimizing boundary tangent cones and single
$C^{1,\theta}$ graphs on their compact regular patches
\cite[Chapter~7]{SimonGMT}. We develop the uniform estimates needed
on those patches, following the regularity strategy of
\cite[Section~2.3]{BiZhu2026} and the relative-forcing analysis of
\cite[Proposition~2.34 and Lemma~2.35]{BiZhu2026v2}.

\begin{lemma}[Simultaneous regular neighborhoods]
\label{lem:rpi:simultaneous-regularity}
Fix a regular point $p$ of one $\Sigma_t$ and a coordinate plane
parallel to its tangent plane there. Given a sufficiently small slope
bound, there are nested coordinate neighborhoods $U\Subset V$ of $p$
such that every frontier meeting $U$ is a single graph in $V$ over
that plane, with the prescribed slope bound and a uniform
$C^{1,\theta}$ estimate. On smaller neighborhoods, every fixed
finite order of one-sided regularity of the graphs, harmonic velocities,
and conformal factors is uniform for $0\le t\le T$.
These conclusions also hold uniformly for the rescaled time families
on each fixed compact regular patch of a limiting cone.
\end{lemma}

\begin{proof}
The proof has two parts. First, ordering prevents a second frontier
from developing excess beside a flat regular patch. Then the harmonic
equation and the minimal graph equation improve the resulting common
graphical bounds together.

\paragraph{A common graphical neighborhood.}
We use the following compactness consequence of the power
almost-minimizing estimate. A sequence with uniform local perimeter
and two-phase density bounds has a locally $L^1$-convergent
subsequence with locally convergent perimeter measures. If the limit
has a regular point, the boundaries converge as single
$C^{1,\theta}$ graphs on a smaller neighborhood of that point,
with uniform bounds. Strict perimeter convergence excludes additional
sheets; $L^1$ convergence alone would not suffice. These are the
almost-minimizer compactness and flat-boundary regularity conclusions
used in \cite[Theorem~2.27]{BiZhu2026v2}.

Fix the reference regular point $p\in\Sigma_t$ and choose coordinates
in which its tangent plane is horizontal. Consider any sequence of
frontiers $\Sigma_{\sigma_j}$ with points
$p_j\in\Sigma_{\sigma_j}$ tending to $p$. After a
subsequence, the filled sets all lie on the same ordered side of
$K_t$. The compactness just stated gives a limiting filled set $K_*$,
and two-phase density puts $p$ on its perimeter support. Ordering
passes to the limit. Every tangent filled set of $K_*$ at $p$ is
therefore ordered with the tangent half-space of $K_t$; its
multiplicity-one minimizing boundary cone is supported in a closed
half-space.

Such a cone is a plane. Indeed, the normal coordinate is nonnegative
on the cone and satisfies
$\Delta_{\rm link}x_\perp=-(k-1)x_\perp$ on its spherical link.
Integration forces that coordinate to vanish. Local area growth and
the singular-set codimension at least seven supply cutoffs with
vanishing tangential $L^1$ gradient, so the integration includes the
singular set. Multiplicity one and the two-phase density bounds give
the nonempty plane with density one. Density-one regularity thus
makes $\partial K_*$ a single $C^{1,\theta}$ graph near $p$, tangent
to the reference sheet. Regular-patch convergence gives the same
fixed graphical neighborhood and a uniform bound for
$\Sigma_{\sigma_j}$ when $j$ is large.

Here is the uniformity conclusion. Fix a small slope threshold and
choose a corresponding bound for the scaled $C^{1,\theta}$ graph
norm in the flat-boundary regularity theorem. If no neighborhoods
$U\Subset V$ worked for the whole time family, choose radii
$r_j\downarrow0$ and frontiers meeting $B_{r_j^2}(p)$ at points
$p_j$ that fail this graph bound on the radius-$r_j$ cylinder
centered at $p_j$. The preceding subsequence instead converges on
a fixed graphical neighborhood of $p$. Its slopes at $p_j$ tend to
the horizontal slope, and their oscillation on a radius-$r_j$
cylinder tends to zero under the uniform $C^{1,\theta}$ bound.
A fixed shrink of the cylinders allows their centers to move from
$p$ to $p_j$. Restricting the fixed graphs then contradicts the
chosen failures. This proves the common graphical neighborhood.

For a rescaled family, repeat this contradiction in a fixed regular
cone chart. The reference graphs converge to its smooth sheet, the
rescaled smooth background metrics have uniform bounds, and the
power errors tend to zero. The same ordered-limit and density-one
argument gives fixed inner and outer graphical neighborhoods,
uniformly in the rescaling index. A finite cover gives the assertion
on a compact regular cone patch. None of these arguments uses
separation for frontiers in different conformal metrics.

\paragraph{One-sided derivatives and the first graph equation.}
Choose three nested neighborhoods
$V_0\Subset V_1\Subset V_2$ within this common graphical region.
For any time $\sigma$, if $\Sigma_\sigma$ meets $V_1$, its uniform
$C^{1,\theta}$ graph and the zero Dirichlet trace of $v_\sigma$ give
one-sided boundary estimates
\[
 \|v_\sigma\|_{C^{1,\alpha}(\overline{E_\sigma}\cap V_0)}\le C,
 \qquad 0<\alpha<\theta.
\]
Here and below norms on graph domains use their uniformly controlled
charts. If $\Sigma_\sigma$ misses $V_1$, its distance from $V_0$ is
bounded below. On each relevant component $v_\sigma$ is then harmonic
on a whole neighborhood or identically zero, and interior estimates
give the same bound. It would not suffice merely to know that the
frontier misses $V_0$. Nesting gives $E_t\subset E_\sigma$ for
$\sigma\le t$, so all these exterior estimates can be integrated
on the common later exterior.

For the first stationarity argument, flatten the graphical boundary and
reflect the negative exterior function $v_\sigma$ oddly. Its zero trace
implies that tangential first derivatives vanish at the graph, and odd
reflection matches the normal first derivative. This gives a uniformly
controlled $C^{1,\alpha}$ extension, nonnegative on the interior side,
which dominates the original zero extension. Where a frontier is
absent, keep the harmonic or zero function. On a smaller common
neighborhood, fixed finite charts and a partition of unity preserve
agreement on the exterior and domination on the interior. Integrating
these extensions gives
\[
 \widetilde u_t\ge u_t,\qquad
 \widetilde u_t=u_t\quad\hbox{on }\overline{E_t},\qquad
 \|\widetilde u_t\|_{C^{1,\alpha}}\le C.
\]
Shrinking the neighborhood preserves a uniform positive lower bound.
Thus $\widetilde g_t=\widetilde u_t^{4/d}g_0$, $d=n-2$, is a
uniformly controlled $C^{1,\alpha}$ extension of the one-sided metric.

These integrals require only measurable, not continuous, time dependence.
The original harmonic velocities are pointwise monotone in time,
interior derivatives are limits of measurable difference quotients,
and nested graph heights are measurable. Reflection in the graphs and
the finite chart constructions preserve this measurability. Equivalently,
one can first integrate the derivatives directly on the common later
exterior by dominated convergence and use extensions only for the graph
equation.

By Lemma~\ref{lem:rpi:flow-construction}, each $K_t$ is locally
perimeter minimizing in $g_t$, including for variations across $K_0$.
Since $\widetilde g_t\ge g_t$ and the two area functionals agree on
$\Sigma_t$, it is also locally minimizing in $\widetilde g_t$.
The graph is stationary in this
extension metric. The minimal graph equation and uniform elliptic
regularity improve the graphs to $C^{2,\alpha}$ on smaller disks.
The constants depend on the graph and metric bounds. The common
local minimizing comparison makes the equation valid on the whole
regular graph disk, including where that disk meets the initial
conductor. At time zero the metric is the smooth reference metric
and the same local minimizing comparison holds. Thus the estimate is uniform for
the whole time family, including zero.

\paragraph{Closing the bootstrap.}
We now have a common $C^{2,\alpha}$ graph bound, including at time
zero. Boundary estimates on these improved domains give uniform
$C^{2,\alpha}$ estimates for all harmonic velocities. Integrate them
on each common later exterior to improve all conformal factors, and
apply the minimal graph equation again. Repetition on nested
neighborhoods gives every prescribed finite order. At these later
stages use a bounded extension operator for uniformly regular graph
domains. Higher-order extensions need not dominate the metric:
stationarity is already established, and agreement of the boundary
value and first derivatives gives the same minimal graph equation.
In particular, no uniform positive Hopf lower bound is needed to
control the sign of higher-order extensions.

For each fixed order only finitely many nested neighborhoods are used.
On a compact regular cone patch the rescaled smooth background metrics
have uniform bounds of every such order, while the lower bound for
$u_t$ and the bound for $v_t$ are unchanged. All the estimates are
therefore uniform in the blow-up index. Together with the original
$C^1$ graphical convergence, interpolation gives convergence in every
lower fixed differentiability order. In particular the one-sided
working metrics converge to their constant contact-point values with
the derivatives needed in the Jacobi expansion.
\end{proof}

The preceding lemma controls derivatives absolutely. Separation also
requires a relative estimate: the error made by transferring a normal
derivative between nearby graphs must be small compared with that
derivative, even when the harmonic function itself tends to zero.

\begin{lemma}[Relative harmonic derivatives]
\label{lem:rpi:relative-harmonic}
Consider uniformly $C^{2,\alpha}$ graphical domains in a fixed
cylinder, with uniformly elliptic $C^{2,\alpha}$ metric coefficients.
Let $z_i>0$ be harmonic on the exterior side, with zero boundary trace,
and let $A_i$ be interior points whose distances from the graph and
the lateral boundary are bounded below by a fixed positive fraction of
the cylinder radius. These are the interior corkscrew points used in
the estimate. Put
$a_i=z_i(A_i)$. On nested central cylinders,
\[
 \|z_i\|_{C^{2,\alpha}}\le C a_i,
 \qquad c a_i\le\partial_\nu z_i\le C a_i
\]
on the original boundary, where $\nu$ points into the exterior.
If a second graph lies in the closure of that exterior and converges
in $C^1$ to the first, then on the second graph the derivative in its
corresponding exterior-pointing unit normal also lies between
$c a_i$ and $C a_i$, with adjusted fixed constants, for large $i$.
All assertions are relative to $a_i$; no positive lower bound for
$a_i$ is required.
\end{lemma}

\begin{proof}
Normalize $z_i$ by $a_i$. Interior Harnack chains and the boundary
Carleson estimate bound this normalized function on a smaller cylinder.
Boundary Schauder estimates then give the stated $C^{2,\alpha}$ bound.
An interior Harnack chain supplies a fixed positive comparison value
on a sphere in a uniform interior tangent ball. The tangent-ball
barrier gives the positive lower bound for the normal derivative on
the central boundary patch; the upper bound follows from the derivative
estimate. These are the standard local boundary estimates for uniformly
regular elliptic Dirichlet problems; see \cite{GilbargTrudinger}.

Join corresponding points of the two graphs by the common graph
coordinate segment, which lies on the exterior side of the first
graph. The Hessian bound and convergence of the normals show that
the error in transferring the derivative is at most
\[
 C a_i\bigl(\hbox{height difference}+\hbox{normal difference}\bigr).
\]
For large $i$ this is less than half the first positive lower bound.
The same factor $a_i$ multiplies the bound and its error, even if
$a_i\to0$. This proves the relative assertion.
\end{proof}

\paragraph{Applying the relative estimate to two flow times.}
Fix $0\le s<t\le T$, and suppose a simultaneous blow-up at a
contact point has a common regular cone patch. We apply the lemma to
the positive functions whose integral measures the change from $u_s$
to $u_t$. Nesting
traps every intermediate filled set between the endpoints. Uniform
almost-minimizing compactness and small-excess regularity give uniform
single-sheet $C^{1,\theta}$ graphs for all intermediate times: an
offending sequence of times would still have the squeezed
multiplicity-one set and perimeter limit, contradicting small-excess
regularity. In a common graph direction the intermediate heights lie
between the endpoint heights. They converge uniformly in $C^0$ and,
by interpolation with the uniform graphical bound, in $C^1$ on a
smaller patch. Nesting alone would not imply the latter statement.
Lemma~\ref{lem:rpi:simultaneous-regularity} supplies the higher-order
bounds for the whole family.

The harmonic functions used in
Lemma~\ref{lem:rpi:relative-harmonic} must be
\begin{equation}\label{eq:rpi:relative-integrands}
 z_\sigma=-\frac{v_\sigma}{u_s}\quad\hbox{on }E_\sigma,
 \qquad s\le\sigma\le t.
\end{equation}
Both numerator and denominator are $g_0$-harmonic on that domain, and
direct calculation gives
\[
 \Delta_{g_s}\left(\frac{f}{u_s}\right)
 =u_s^{-(n+2)/d}
       \left(\Delta_{g_0}f-\frac{f}{u_s}\Delta_{g_0}u_s\right).
\]
Thus $z_\sigma$ is positive and $g_s$-harmonic on its connected
exterior and has zero trace on its own frontier. The denominator is
fixed at time $s$, not at time $\sigma$.

Choose a common corkscrew point $A_i$ in the later exterior, above
all intermediate rescaled graphs. Put
$a_{i,\sigma}=z_\sigma(A_i)>0$. The preceding lemma and the uniform
family bounds give relative estimates independent of $i$ and $\sigma$
on the fixed patch. Transfer each normal derivative to the later
graph. The error is
\[
 C a_{i,\sigma}
  \bigl(\hbox{height difference}+\hbox{normal difference}\bigr)
       =o(1)a_{i,\sigma},
\]
uniformly in $\sigma$. Hence that derivative is between
$c a_{i,\sigma}$ and $C a_{i,\sigma}$. Set
\[
 \lambda_i=\int_s^t a_{i,\sigma}\,d\sigma>0.
\]
The normalization is measurable and bounded in time. Since its
denominator is fixed, the exact identity is
\begin{equation}\label{eq:rpi:relative-integral}
 1-\frac{u_t}{u_s}=\int_s^t z_\sigma\,d\sigma.
\end{equation}
The relative derivative estimate gives an integrable majorant for
exterior difference quotients at the later graph. Dominated convergence
therefore permits differentiation there and proves
\begin{equation}\label{eq:rpi:relative-forcing}
 -C\lambda_i\le
    \partial_{\nu_t}^{g_s}(u_t/u_s)
      \le-c\lambda_i.
\end{equation}
The ratio $u_t/u_s$ is uniformly positive, so the extra positive factor
in the conformal mean-curvature formula preserves these bounds. Only
exterior traces are differentiated; no derivative across the zero
extension of a velocity is taken, and no time derivative at a jump
time is required.

The same estimate applies to the other deformation used below: a
fixed harmonic conformal factor $u$, normalized to have boundary value
one. Its positive harmonic function is $z=1-u$.
The strict inequality $u<1$ on the connected exterior is essential.
Lemma~\ref{lem:rpi:relative-harmonic} then transfers its derivative
from the original frontier to the nearby candidate graph. All constants
may depend on the fixed regular cone patch, finite time interval, and
fixed perturbation parameter. No uniform estimate through a cone
singularity or as the perturbation parameter vanishes is asserted.

\section{Singular separation from a single distance median}
\label{sec:rpi:median-separation}

The next proposition excludes singular contact for an ordered pair
whose regular parts do not meet. The applications below will verify
that regular comparison applies. The idea is to rescale at a possible
singular contact and divide every graph separation by the same number: the median distance from the earlier support to the later
one. This produces one positive supersolution on the regular part of
the common tangent cone. Suitable record scales make its infimum
decay to zero near the vertex. After a bounded truncation and removal
of the singular set by capacity cutoffs, the minimizing-hypersurface
mean-value inequality contradicts that decay.

This follows the singular-separation approach of
\cite[Proposition~2.36]{BiZhu2026}, using the relative forcing,
half-volume normalization, record scales, and truncation developed in
\cite[Proposition~2.34 and Lemma~2.35]{BiZhu2026v2}. We keep the
median independent of the regular-patch exhaustion. That distinction
is essential: a positive Jacobi supersolution alone is possible on a
singular cone---for example, $r^{-2}$ on the seven-dimensional Simons
cone. The contradiction will use the additional small-ball infimum
estimate.

\begin{proposition}[Separation of the ordered supports]
\label{prop:rpi:singular-separation}
Let $K\subset K'$ be canonical closed filled sets in a smooth
$n$-manifold, $n\ge8$, with multiplicity-one perimeter supports
$\Sigma$ and $\Sigma'$. Put $k=n-1$. On a neighborhood of their
possible contact set, let $g_{\rm ref}$ be a smooth metric and assume
that there are $\Lambda<\infty$, $\alpha>0$, and $r_0>0$ such that,
for $E=K,K'$,
\[
 P_{g_{\rm ref}}(E;B_r(x))
 \le P_{g_{\rm ref}}(L;B_r(x))+\Lambda r^{k+\alpha}
 \quad\text{if }L\triangle E\Subset B_r(x),\quad 0<r<r_0.
\]
The balls lie in that neighborhood. Assume that every intersection
point is singular for both supports.

Let $g=a^{4/(n-2)}g_{\rm ref}$, with $a>0$ continuous, be the
metric used to measure graph heights. At a possible contact $p$,
choose coordinates with $g(p)=\delta$ and rescale by
$\eta_i(y)=p+r_i y$, $r_i\downarrow0$. The power comparison gives
simultaneous minimizing-cone subsequences with local perimeter
convergence; the two limiting cones coincide, as shown below.
For every such subsequence with common cone $C$, assume the following
on each compact regular patch, after shrinking a larger regular patch:
\begin{enumerate}
\item The rescaled supports converge as single smooth graphs, and
      extensions of $g_i=r_i^{-2}\eta_i^*g$ converge to $\delta$
      in $C^2$. The graph and metric derivatives needed below have
      uniform local bounds.
\item The later support is the $g_i$-normal graph of a function
      $\zeta_i>0$ over the earlier support, and
\begin{equation}\label{eq:rpi:separation-interface}
 L_i\zeta_i=F_i\le0,\qquad
 L_i\longrightarrow J_C:=\Delta_C+|A_C|^2.
\end{equation}
      In fixed regular-patch coordinates, the second-order operators
      are uniformly elliptic, their coefficients have uniform
      $C^{0,\theta}$ bounds for some $\theta>0$, and converge locally
      to those of $J_C$.
\item On a slightly larger patch there are numbers $\lambda_i>0$
      and constants $c,C>0$, independent of $i$, such that
\[
 -C\lambda_i\le F_i\le0,\qquad
 F_i\le-c\lambda_i\quad\text{on an inner ball}.
\]
      The numbers and constants may depend on the fixed patch, and
      $\lambda_i$ may tend to zero.
\end{enumerate}
Then $\Sigma\cap\Sigma'=\varnothing$.
\end{proposition}

\begin{proof}
Suppose the supports meet at the origin, and put $k=n-1\ge7$.

\paragraph{A common cone with connected regular part.}
The power almost-minimizing comparison gives local strict perimeter
compactness and almost-monotonicity of density. Hence every simultaneous
blow-up has a subsequence converging to multiplicity-one minimizing
boundary cones for $g_{\rm ref}(p)$. They remain nested. Since
$g(p)$ is a constant scalar multiple of $g_{\rm ref}(p)$, both
cones are stationary also in the normalized Euclidean metric.
Their supports are connected, because radial segments join all points
to the vertex; their singular sets have dimension at most $k-7$.
The same-metric strong maximum principle of
\cite[Theorem~1.1 and Remark~(3)]{Wickramasekera2014Maximum} identifies
the two supports, since both contain the vertex. Multiplicity one
identifies their perimeter measures. This works for every simultaneous
blow-up, not only the sequence selected below.

The regular part of the common cone $C$ is connected. To see this,
restrict its varifold to a regular component. Its stationarity extends
across the singular set by the local area bound and cutoffs with
vanishing tangential $L^1$ gradient, since
$\mathcal H^{k-1}(\operatorname{Sing}C)=0$. Each regular component
is dilation invariant: the radial dilation orbit is a path in that
component. Its closure therefore contains the vertex. Distinct component
closures intersect only in the singular set, because a regular point
has a connected smooth-sheet neighborhood. This would give intersecting
stationary integral varifolds whose intersection has zero
$\mathcal H^{k-1}$ measure, contrary to
\cite[Theorem~3.2 and the proof of Theorem~1.1, pp.~8--9]
{Wickramasekera2014Maximum}. These are same-metric statements and do
not assume the different-metric conclusion being proved.

\paragraph{A median chosen before the regular patches.}
Choose a fixed smooth coordinate chart normalized so that the earlier
working metric at the contact point is Euclidean. Let
$D(y)=\dist_{\R^n}(y,\Sigma')$, let $\mu_r$ denote the perimeter
measure of $\Sigma\cap B_r$, and define the fixed median
\begin{equation}\label{eq:rpi:separation-median}
 q(r)=\inf\left\{a\ge0:
   \mu_r(\{D\le a\})\ge\tfrac12\mu_r(B_r)\right\}.
\end{equation}
Atoms are allowed. Both $\{D\le q(r)\}$ and $\{D\ge q(r)\}$
have at least half the measure. The contact set has zero perimeter
measure, since there is no regular contact and the singular sets have
zero perimeter measure. Thus $q(r)>0$ for small positive $r$.
Also $q(r)\le r$, since the origin belongs to $\Sigma'$.

In fact $q(r)/r\to0$. If not, choose a contradicting sequence and
then a simultaneous minimizing-cone subsequence. Compactness and the
density bounds give local Hausdorff convergence of both rescaled
supports to the same cone. On $B_1$, distance to the second support
therefore tends uniformly to zero, contradicting the assumed lower
bound for the rescaled medians.

Put $f(r)=q(r)/r$. Given any $R_i\downarrow0$, choose $0<r_i\le R_i$
such that
\[
 f(r_i)\ge\tfrac12\sup_{0<s\le R_i}f(s).
\]
The supremum is finite and positive; neither continuity of the median
nor attainment of the supremum is required. For every $0<\rho\le1$,
this choice gives
\begin{equation}\label{eq:rpi:median-record}
 \frac{q(\rho r_i)}{q(r_i)}\le2\rho.
\end{equation}
Pass to a simultaneous cone subsequence at these radii, and put
\[
 a_i=\frac{q(r_i)}{r_i}>0,
 \qquad \Sigma_i=r_i^{-1}\Sigma,
 \qquad \Sigma'_i=r_i^{-1}\Sigma'.
\]
Every graph separation on every patch will be divided by this same
$a_i$.

\paragraph{From ambient distance to normal height.}
The record scales and the normalizing numbers $a_i$ are now fixed.
We next relate their ambient-distance definition to the graph variable
in \eqref{eq:rpi:separation-interface}. On a compact regular core of $C$,
choose a slightly larger regular neighborhood. For large $i$,
$\Sigma'_i$ is the $g_i$-normal exponential graph of $\zeta_i>0$
over $\Sigma_i$, where $g_i$ is the rescaled earlier working metric.
Their relative slopes tend to zero. For a point in the smaller core,
its nearest point on $\Sigma'_i$ lies in the larger neighborhood:
the local normal candidate has distance tending to zero, while the
complement of that neighborhood stays a fixed positive distance away.
The Lipschitz-graph distance estimate in coordinates flattening
$\Sigma_i$ gives
\begin{equation}\label{eq:rpi:distance-height}
 \dist_{\R^n}(y,\Sigma'_i)=(1+o(1))\zeta_i(y)
\end{equation}
uniformly on the smaller core. The normal exponential segment has
Euclidean length $(1+o(1))\zeta_i$ and direction approaching the
Euclidean normal, because the contact metric was normalized and the
rescaled metrics converge there in $C^1$. Thus the same relative
Lipschitz-graph estimate applies to its endpoint. This remains a
relative estimate even when $\zeta_i$ is much smaller than the
blow-up error; an additive comparison with the cone would not suffice.
Figure~\ref{fig:rpi:median-normalization} distinguishes these two
measurements and the use of one median on all regular patches.

\begin{figure}[tb]
\centering
\begin{tikzpicture}[x=1cm,y=1cm,font=\small,
  every node/.style={inner sep=2pt},
  support/.style={line width=.85pt},
  measure/.style={line width=.75pt},
  graphheight/.style={line width=.9pt,dashed,blue!55!black}]
  \node[font=\small\bfseries] at (3.35,4.15)
    {One median for the whole ball};
  \node[font=\small\bfseries] at (10.65,4.15)
    {Two measurements on a regular patch};

  \node[align=center] at (3.35,3.55)
    {$D_i(y)=\dist_{\mathbb R^n}(y,\Sigma'_i)$};
  \node[align=center] at (3.35,2.95)
    {$\displaystyle a_i=\frac{q(r_i)}{r_i}$};
  \node[align=center] at (3.35,2.03)
    {$\mu_i(\{D_i\le a_i\})\ge\tfrac12\mu_i(B_1)$\\[4pt]
     $\mu_i(\{D_i\ge a_i\})\ge\tfrac12\mu_i(B_1)$};
  \draw[-{Stealth[length=2mm]},black!55] (3.35,1.48) -- (3.35,.98);
  \node[align=center] at (3.35,.47)
    {$\displaystyle W_i=\frac{\zeta_i}{a_i}$\\[4pt]
     on every compact regular patch};

  \draw[black!20] (6.9,-.1) -- (6.9,3.85);

  \draw[support] (7.55,1.02)
    .. controls (8.55,.74) and (9.8,.95) .. (10.65,.95)
    .. controls (11.45,.95) and (12.65,.85) .. (13.65,1.1);
  \node[anchor=west] at (13.7,1.1) {$\Sigma_i$};
  \draw[support,blue!55!black] (7.55,2.204) -- (13.65,3.18);
  \node[anchor=west,blue!55!black] at (13.7,3.18) {$\Sigma'_i$};

  \coordinate (y) at (10.65,.95);
  \coordinate (P) at (10.37699,2.65632);
  \coordinate (X) at (10.88,2.7368);
  \draw[measure] (y) -- (P)
    node[pos=.48,left=5pt] {$D_i(y)$};
  \draw[graphheight] (y)
    .. controls (10.82,1.48) and (10.94,2.14) .. (X)
    node[pos=.51,right=5pt] {$\zeta_i(y)$};
  \fill (y) circle (1.7pt);
  \node[below=5pt] at (y) {$y$};
  \fill (P) circle (1.5pt);
  \fill[blue!55!black] (X) circle (1.5pt);
  \node[anchor=south east] at (10.22,2.96) {nearest point};
  \draw[black!55,line width=.4pt] (10.17,2.95) -- (P);
  \node[anchor=south west,blue!55!black] at (11.2,3.35)
    {normal-graph point};
  \draw[blue!55!black,line width=.4pt] (11.26,3.33) -- (X);
  \node[align=center] at (10.65,.12)
    {$D_i=(1+o(1))\zeta_i$ on compact regular cores};
\end{tikzpicture}
\caption{A single distance median and its use on regular patches.
The curves are schematic regular cross-sections of the rescaled ordered
supports; no graph through the singular set is depicted.
The straight segment measures the ambient Euclidean nearest-support
distance $D_i(y)$, whereas the dashed segment is the normal exponential
segment of length $\zeta_i(y)$ in the rescaled earlier working metric
$g_i$. For $T_i(x)=x/r_i$, the measure
$\mu_i=r_i^{-k}(T_i)_\#\mu_{r_i}$ is the rescaling of the same
perimeter measure used to define $q(r_i)$, on all of $\Sigma_i\cap B_1$.
The number $a_i=q(r_i)/r_i$ is fixed using that whole measure before
regular patches are chosen. Both median
inequalities allow atoms, and every patch uses the same $a_i$.
The displayed relative comparison is uniform on each compact regular
core as $i\to\infty$.}
\label{fig:rpi:median-normalization}
\end{figure}


Choose a compact regular core inside $C\cap B_1$ whose omitted
area is less than one eighth of its total area, and enlarge it
slightly while keeping away from the singular set and sphere.
Such cores exist because those two sets have zero $k$-measure on
the cone. Graphical perimeter convergence makes the omitted area
less than one quarter for large $i$. Each median set therefore
retains positive measure on this fixed core.

\paragraph{Local bounds with the same normalization on every patch.}
Set $W_i=\zeta_i/a_i$. The lower median set and
\eqref{eq:rpi:distance-height} provide a point with $W_i\le2$ in
one of finitely many regular core patches. Choose that patch along
a subsequence. Weak Harnack for the nonnegative supersolution
$L_iW_i\le0$ gives a uniform local $L^p$ bound for some $p>0$,
before any bound for $F_i/a_i$ has been proved. To handle the
zero-order term, write
\[
 L_i=A_i^{ab}\partial_{ab}+B_i^a\partial_a+c_i
\]
and replace $c_i$ by $\min\{c_i,0\}$, obtaining $\widehat L_i$.
Since $W_i\ge0$, we still have $\widehat L_iW_i\le0$.
The usual weak Harnack inequality with nonpositive zero-order
coefficient applies with uniform constants on every fixed patch.

An $L^p$ bound on an overlap of fixed positive measure supplies a
bounded-value point for the next patch. Weak Harnack there propagates
the bound. Connectedness of $\operatorname{Reg}C$ gives finite
chains from the anchor patch to every fixed compact regular set.
Consequently
\begin{equation}\label{eq:rpi:median-local-lp}
 \|W_i\|_{L^p(K)}\le C_K
 \quad\hbox{for every }K\Subset\operatorname{Reg}C.
\end{equation}
The constants may grow with $K$. No uniform Harnack chain through
the singular set has been asserted.

\paragraph{Controlling the forcing before passing to a limit.}
The local $L^p$ bounds have been obtained without dividing a
right-side estimate by the possibly small number $a_i$. They now
let us control that quotient. In a ball
where $F_i\le-c\lambda_i$, solve
\[
 \widehat L_i\phi_i=-1,\qquad \phi_i|_{\partial B}=0.
\]
Uniform ellipticity, coefficient bounds, and the maximum principle
give $\phi_i\ge c_0>0$ on a smaller ball. Since
$\widehat L_i\zeta_i\le F_i$ and $\zeta_i\ge0$ on the boundary,
comparison gives
\[
 \zeta_i\ge c\lambda_i\phi_i.
\]
Combining this with \eqref{eq:rpi:median-local-lp} yields
$\lambda_i/a_i\le C_K$. The assumed upper bound for $|F_i|$
and a finite cover now give
\begin{equation}\label{eq:rpi:median-forcing-bound}
 \|F_i/a_i\|_{L^\infty(K)}\le C_K.
\end{equation}
The order is important: first weak Harnack, then the positive
Dirichlet comparison, and only then bounded-right-side elliptic
estimates.

The weak Harnack exponent need not exceed one. The local boundedness
estimate for the now bounded-right-side equation bounds the supremum
on an inner ball by the $L^2$ norm on a larger ball and the right-side
bound. Interpolate the $L^2$ norm between the bounded $L^p$ norm
and the larger-ball supremum, apply Young's inequality, and iterate
on nested radii to absorb that supremum. For each $i$ the preliminary
finite suprema exist by local regularity. This gives uniform local
$L^\infty$ bounds; interior elliptic estimates then give uniform
$W^{2,q}$ bounds for every finite $q$. One countable exhaustion and
one diagonal subsequence give a single function $W$, with local
$C^1$ convergence on all of $\operatorname{Reg}C$, satisfying
\begin{equation}\label{eq:rpi:median-limit}
 W\ge0,
 \qquad W\in W^{2,q}_{\rm loc}(\operatorname{Reg}C)
          \quad(q<\infty),
 \qquad (\Delta_C+|A_C|^2)W\le0.
\end{equation}
The upper median set also retains positive measure on the fixed
regular core. Equation~\eqref{eq:rpi:distance-height} gives points
there with $W_i\ge1/2$. Local uniform convergence makes $W$
nonzero. Since $\Delta_CW\le-|A_C|^2W\le0$, the strong maximum
principle and connectedness make $W>0$ throughout the regular part.
No separately normalized patch limits are used.

\paragraph{The infimum near the vertex.}
We have obtained one positive limit on the entire regular part. The
record-scale choice now gives information on its behavior near the
vertex without changing its normalization or taking new patchwise
limits.

Fix $0<\rho\le1$. The lower median set in $\Sigma_i\cap B_\rho$
has half its area and, by \eqref{eq:rpi:median-record}, normalized
ambient distance at most $2\rho$. Choose a compact regular core
in $C\cap B_\rho$ omitting less than one quarter of its area.
Perimeter convergence and \eqref{eq:rpi:distance-height} give a
point in this core with $W_i\le3\rho$ for large $i$. The already
chosen local uniform convergence applies to this core. Thus the
same limit satisfies
\begin{equation}\label{eq:rpi:median-small-infimum}
 \inf_{\operatorname{Reg}C\cap B_\rho}W\le3\rho.
\end{equation}
It is enough to use a countable sequence of $\rho\downarrow0$.
The constant $3$ does not depend on the core or on $\rho$: for
each fixed radius the relative error in
\eqref{eq:rpi:distance-height} tends to zero before this limit is
taken. Therefore the infimum of $W$ in every ball about the vertex
is zero. Its essential infimum there is also zero, since every
small regular value has a neighborhood of positive perimeter measure
with comparably small values.

\paragraph{A bounded supersolution across the singular set.}
To apply a mean-value inequality on the whole minimizing cone, we
need a Sobolev supersolution across its singular set. First discard
the nonnegative Jacobi potential, obtaining $\Delta_CW\le0$, and
then set $h=\min\{W,1\}$. Approximation by increasing concave truncations
shows that $h$ is a bounded positive weak \emph{Laplace}
supersolution on the regular part. It need not be a Jacobi
supersolution. For a regular-part cutoff $\eta$, use the nonnegative
test $\eta^2(1-h)$ to obtain
\[
 \int\eta^2|\nabla_C h|^2
 \le2\int\eta(1-h)|\nabla_C h||\nabla_C\eta|,
\]
and hence
\begin{equation}\label{eq:rpi:truncation-energy}
 \int\eta^2|\nabla_C h|^2
       \le4\int|\nabla_C\eta|^2.
\end{equation}

The singular set has zero tangential two-capacity on compact regions.
Indeed, cover the relevant compact singular set by finitely many balls
$B_{\rho_j}(x_j)$, with centers on that set, maximum radius tending
to zero, and $\sum_j\rho_j^{k-2}\to0$. This is possible because
its dimension is at most $k-7$. Let $\theta_j$ vanish on the ball,
equal one off its double, and satisfy $|D\theta_j|\le C/\rho_j$.
Put $\chi=\min_j\theta_j$. Almost everywhere its gradient is one
of the active gradients, so local area growth gives
\begin{equation}\label{eq:rpi:cone-capacity}
 \int_C|\nabla_C\chi|^2
 \le\sum_j\int_{C\cap B_{2\rho_j}(x_j)}|\nabla_C\theta_j|^2
 \le C\sum_j\rho_j^{k-2}\longrightarrow0.
\end{equation}
There is no overlap-count factor. These cutoffs vanish near the
whole relevant singular set and converge to one on the regular part.

Use $\eta=\psi\chi$ in \eqref{eq:rpi:truncation-energy}, where
$\psi$ is a fixed compact spatial cutoff. Fatou's lemma gives
locally finite Dirichlet energy for $h$ across the singular set.
Moreover, $\psi h\chi\to\psi h$ in the local surface
$W^{1,2}$ norm: the term $h\nabla_C\chi$ tends to zero by
boundedness and \eqref{eq:rpi:cone-capacity}, while the omitted
$\nabla_C h$ term tends to zero by absolute continuity of its
now finite energy. Thus $h$ genuinely extends as a local Sobolev
function on the whole minimizing boundary; its values on the
measure-zero singular set are immaterial.

For a nonnegative ambient smooth compactly supported test function
$\varphi$, insert $\varphi\chi$ into the regular-part supersolution
inequality. The cutoff error is bounded by
\[
 \|\varphi\|_\infty
 \|\nabla_C h\|_{L^2(\supp\varphi)}
 \|\nabla_C\chi\|_2\longrightarrow0.
\]
The other term converges by dominated convergence. Hence
\[
 \int_C\langle\nabla_C h,\nabla_C\varphi\rangle\ge0
\]
for tests meeting the singular set as well.

\paragraph{The mean-value contradiction.}
The bounded function $h$ is now a positive weak Laplace supersolution
on the whole minimizing boundary, with locally finite energy and zero
essential infimum in every ball about the vertex. These are the
conditions needed for the minimizing-boundary weak Harnack theorem
of Bombieri--Giusti \cite{BombieriGiusti1972Harnack}; the precise
form needed here is stated in
\cite[Theorem~3.1, p.~11]{DeSilvaJerison2011Gradient}.
For an area-minimizing hypersurface in $B_R$ and a positive weak
Laplace supersolution, a normalized $L^p$ mean on $B_r$ is bounded
by its infimum on $B_r$ when $r\le\beta R$ and
$0<p<k/(k-2)$. The constants depend only on dimension and $p$.
This is an extrinsic estimate on the minimizing boundary, not an
assertion that regular-patch Harnack chains have uniform constants
through cone singularities.

Apply it to $h+\epsilon$ on the cone, with fixed concentric balls
and, for example, $p=1$. Choose the smaller ball to contain a regular
open subset on which $h>0$; the cone and the same $h$ are defined
on any required larger ball. Its conclusion is
\[
 \frac{1}{\mathcal H^k(C\cap B_r)}
       \int_{C\cap B_r}(h+\epsilon)
 \le C\operatorname*{ess\,inf}_{C\cap B_r}(h+\epsilon).
\]
The essential infimum of $h$ is zero by
\eqref{eq:rpi:median-small-infimum}. Sending $\epsilon\downarrow0$
would make the positive integral on the left zero, a contradiction.
This proves separation. No growth selection at infinity or limiting
truncation level is needed.
\end{proof}

\subsection{The graph equation for the two deformations}
\label{sec:rpi:separation-application}

It remains to verify the separation proposition for the flow and for
the prescribed-curvature enclosure. Both use the same graph equation;
the difference is a small zero-order term in the latter problem. We
derive the equation using the end-pointing normal throughout.

Normalize the earlier working metric $g$ at the contact point to be Euclidean, write $g_i$
for its rescaling at radius $r_i$, and use $g_i$-unit normals and
$g_i$-normal exponential graphs throughout. For this calculation assume uniform $C^{3,\alpha}$ bounds for the
extended metrics and graphs on a larger regular cone patch.
Lemma~\ref{lem:rpi:simultaneous-regularity} supplies them for flow
slices; the candidate construction below will supply them in the
second application.
Continuity of the original metric gives uniform convergence to its
contact-point value. Apply bounded graph-domain extension operators
to its difference from that constant metric; these preserve uniform
convergence. Interpolation then gives $g_i\to\delta$ in $C^2$
on a smaller patch. The graph bounds similarly give the $C^2$
convergence needed below. A mere $C^{1,\alpha}$ metric bound would
not justify convergence of the Ricci term.

The linearization of $H=\divg\nu$ at the earlier minimal graph is
$-J_i$, where
\[
 J_i=\Delta_{\Sigma_i,g_i}+|A_i|_{g_i}^2
                          +\Ric_{g_i}(\nu_i,\nu_i).
\]
Integrating the derivative of the mean-curvature graph operator from
height zero to $\zeta_i$ writes its nonlinear remainder as
\[
 Q_i\zeta_i
   =Q_i^{ab}\partial_{ab}\zeta_i
      +Q_i^a\partial_a\zeta_i+Q_i^0\zeta_i,
\]
with coefficients tending locally to zero. Let the later metric be
$g'_i=U_i^{4/d}g_i$, and let its prescribed end-oriented curvature
be $f_i(y)=r_i f(\eta_i(y))$, where $\eta_i(y)=p+r_i y$.
The conformal formula gives the exact equation
\begin{equation}\label{eq:rpi:separation-graph-equation}
 (J_i-Q_i)\zeta_i=F_i-U_i^{2/d}f_i,
 \qquad F_i=\frac{2k}{d}U_i^{-1}
                         \partial_{\nu'_i}^{g_i}U_i.
\end{equation}
For flow slices $f=0$ and $U_i$ is the rescaled ratio $u_t/u_s$.
Equations~\eqref{eq:rpi:relative-integrands}--
\eqref{eq:rpi:relative-forcing} give the required two-sided relative
bound for $F_i$.

\subsection{Separation of flow slices}
\label{sec:rpi:flow-slice-separation}

Fix $0\le s<t$. We prove that the two closed supports are disjoint.
Suppose $p\in\Sigma_s\cap\Sigma_t$ is regular
for at least one frontier. The simultaneous regular-neighborhood
lemma makes both graphs regular there. Nesting puts $p$ on every
intermediate frontier: an interior point of an intermediate conductor
would be interior to $K_t$. These ordered graphs have a common
tangent plane and the same end-oriented normal. Their velocities
vanish at $p$, so $u_t(p)=u_s(p)$.

For each $\sigma\in[s,t]$, the function
$z_\sigma=-v_\sigma/u_s$ is positive and $g_s$-harmonic on its
own exterior $E_\sigma$, vanishes at $p$, and has strictly positive
derivative in the common $g_s$-unit normal direction. Uniform
boundary $C^1$ bounds and the positive lower bound for $u_s$ give
an integrable bound for these measurable derivatives. The uniform
$C^2$ common-tangent graph charts provide a single short outward
normal segment in $E_t$, hence in every $E_\sigma$; all exterior
difference quotients can be taken there. Thus, with $q=u_t/u_s$,
Equation~\eqref{eq:rpi:relative-integral} gives
\[
 \partial_\nu q(p)
       =-\int_s^t\partial_\nu z_\sigma(p)\,d\sigma<0.
\]
Extend the later metric from its exterior as a $C^1$ metric. Its
first jet at $p$ is determined by that one-sided metric, and the
earlier metric has the matching one-sided jet. The conformal
mean-curvature formula at this point gives
\[
 H_{\Sigma_s}^{g_t}(p)
 =q(p)^{-2/d}\left(H_{\Sigma_s}^{g_s}(p)
                +\frac{2k}{d}q(p)^{-1}\partial_\nu q(p)\right)<0.
\]
The earlier curvature is zero, including at $s=0$. The later graph
has zero curvature in the same extension metric. Since $K_s\subset K_t$,
the tangent graph comparison gives
$H_{\Sigma_t}^{g_t}(p)\le H_{\Sigma_s}^{g_t}(p)<0$, a contradiction.
Only the common exterior first jets were used; the ratio need not be
constant on the later frontier or harmonic throughout $E_s$.
Thus regular contact is impossible. At any remaining contact, the
power almost-minimizing estimates and density bounds give the common
minimizing-cone limits used in
Proposition~\ref{prop:rpi:singular-separation}. The regular-patch
bootstrap, relative harmonic derivative estimate, and graph equation
verify its remaining hypotheses. The proposition excludes singular
contact as well, proving $\Sigma_s\cap\Sigma_t=\varnothing$.

\subsection{The prescribed-curvature enclosure before detachment}
\label{sec:rpi:perturbed-candidate}

We now build the enclosing frontier needed for the mass--capacity
comparison. The construction applies both to a flow slice and to the
original smooth ambient metric in the approximation argument. In
either case, the initial frontier is locally perimeter minimizing;
its enclosing candidate must be understood before detachment is known.

Let $K$ be a compact canonical closed filled set with nonempty
frontier $\Sigma$ and connected open end exterior $E$. Work in either
of the following settings:
\begin{enumerate}
\item $K=K_s$ is a flow conductor and $h=b^2g_s$ for a fixed
      constant $b>0$;
\item $h$ is smooth through $\Sigma$ and on $E$, and $K$ is
      locally perimeter minimizing in $h$ on a neighborhood of
      $\Sigma$. The exterior has one coordinate end and compact
      complement, with $h_{ab}=a_0^2\delta_{ab}+O_2(r^{-\tau_0})$
      for constants $a_0,\tau_0>0$.
\end{enumerate}
In the second case extend $h$ smoothly a short distance behind the
frontier and, if needed, use the buried-boundary construction already
described. No outermost property of $K$ is required. In both settings
$K$ is locally minimizing in its continuous metric $h$: for a flow
slice this is Lemma~\ref{lem:rpi:flow-construction}, and in the
smooth setting it is the stated hypothesis. The density and minimizing
cone conclusions follow from the same local comparison.

Let $u$ be a positive $h$-harmonic function on $E$, continuous up to
$\Sigma$, with $u=1$ there, $u<1$ in $E$, and $u\to\ell\in(0,1)$
at infinity. Extend it by one on $K$ and set $h'=u^{4/d}h$.
Assume on the coordinate end that
$h'_{ab}=a^2\delta_{ab}+O_2(r^{-\tau})$, where $a,\tau>0$.
This differentiated decay will confine the candidate inside a fixed
coordinate sphere.

The extended perturbation has the H\"older modulus needed for the
weighted perimeter comparison. To see this directly in the flow
case, the conformal Laplacian identity makes
\[
 z=u_s(1-u)
\]
a bounded positive $g_0$-harmonic function on $E_s$, continuous with
zero frontier value and end value
$c=e^{-s}(1-\ell)>0$. Let $q_{K_s}$ be the escape potential for
$g_0$, constructed in the proof of
Lemma~\ref{lem:rpi:flow-construction}. Bounded continuous Dirichlet
uniqueness gives $z=cq_{K_s}$. The already proved H\"older modulus
of $q_{K_s}$ and the positive H\"older factor $u_s$ therefore give
\[
 u=1-\frac{cq_{K_s}}{u_s}
\]
with a H\"older extension across $\Sigma$. For a smooth initial
metric, use the same background capacitary construction with $h$ in
place of $g_0$. Local minimizing density gives its boundary
contraction estimate, and uniqueness gives
$1-u=(1-\ell)q_K$. Thus $h'$ is a positive H\"older metric in
both cases, before the candidate is constructed.

Choose a smooth ambient function $h_0$ that is zero on $K$, positive
on $E$, and integrable in $h'$. A locally finite partition of unity
in $E$, with coefficients decreasing sufficiently rapidly, makes its
zero extension smooth and its integral finite. Its first derivative
vanishes along $\Sigma$; in particular, on compact sets,
$|h_0(x)|\le C\operatorname{dist}(x,\Sigma)$.

For $\delta>0$ chosen below, the candidate $K'$ minimizes
\[
 \mathcal F(L)=P_{h'}(L)
       +\delta\int_{L\setminus K}h_0\,dV_{h'}
\]
among all compact finite-perimeter supersets of $K$ in the ambient
manifold. The added region $L\setminus K$ need not stay a positive
distance from the original frontier. The parameters, forcing, and its derivatives are
fixed here. To fix the sign, a normal variation $\varphi\nu$ of the candidate
frontier has first variation
\[
 D\mathcal F[\varphi\nu]
   =\int_{\partial K'}(H_\nu+\delta h_0)\varphi\,dA_{h'}.
\]
Its prescribed end-oriented curvature is therefore $H=-\delta h_0$.
The function $h_0$ vanishes on the original filled side; its value on
the candidate frontier has not been prescribed to be zero.

\paragraph{Existence, compact support, and a connected end exterior.}
On the asymptotically flat end, $h'$ is smooth and approaches a positive constant multiple
of the Euclidean metric, with the corresponding differentiated decay.
Choose $R_*$ enclosing $K$ so that every coordinate sphere beyond
$R_*$ has strictly positive outward mean curvature in $h'$. This
choice precedes that of $\delta$. First minimize $\mathcal F$ over
all finite-perimeter supersets of $K$ in a coordinate truncation
$B_Q$, including the compact core, with $Q>R_*$.
The competitor $K$ gives a perimeter bound. BV compactness, lower
semicontinuity for the positive continuous perimeter weight, and
$L^1$ convergence of the bounded volume density on $B_Q$ give a
minimizer. If $X$ is the outward $h'$-unit normal to the coordinate
spheres, the divergence formula gives, for almost every $r>R_*$,
\[
 P_{h'}(L\cap B_r)
 \le P_{h'}(L)
       -\int_{L\setminus B_r}\operatorname{div}_{h'}X\,dV_{h'}.
\]
The divergence is positive. Clipping also decreases the nonnegative
volume term, strictly decreasing $\mathcal F$ whenever the removed
tail has positive volume. The minimizer is therefore supported in
$B_{R_*}$, independently of $Q$ and $\delta$. Clipping an arbitrary
compact competitor into this same ball proves global minimality,
with no artificial outer constraint remaining.

To ensure a filled enclosure, choose $\delta>0$ sufficiently small
after fixing $h'$ and $R_*$. Compact-support isoperimetry gives
\[
 \Vol_{h'}(U)^{k/n}\le C_{\rm iso}P_{h'}(U)
\]
for sets $U\subset B_{R_*}$ disjoint from the buried collar. Such
sets have zero trace on a slightly larger smooth truncation, so
this is the absolute inequality. Require
\[
 \delta\|h_0\|_{L^\infty(B_{R_*})}
 C_{\rm iso}\Vol_{h'}(B_{R_*})^{1/n}<1.
\]
In a larger truncation $B_{Q_2}$, decompose $B_{Q_2}\setminus L$
into its finite-perimeter components. Exactly one contains the
connected outer annulus; let $U$ be the union of the others.
Perimeter additivity, with cancellation of the outer sphere, shows
that filling $U$ removes its perimeter. Consequently
\[
 \mathcal F(L\cup U)-\mathcal F(L)
 =-P_{h'}(U)+\delta\int_Uh_0\,dV_{h'}<0
\]
unless $U$ has zero volume. This contradicts minimality. The selected
minimizer is thus filled and has connected end exterior, while
retaining minimality against all finite-perimeter enclosing
competitors. No positive lower bound on $\delta$ is needed.

\paragraph{Local variations across the original frontier.}
The existence argument has produced a global enclosing minimizer.
To obtain its graph equation near possible contact, we remove the
enclosure constraint for local variations. Let $L\triangle K'$ be
compactly supported in a sufficiently small variation ball. Since
$h'=h$ on $K$, the local minimizing property of $K$ gives
\[
 P_{h'}(L\cap K)\ge P_{h'}(K),\qquad
 P_{h'}(L\cup K)\le P_{h'}(L),
\]
where the second inequality is perimeter submodularity and all
comparisons are localized to that ball. This argument uses only local
minimality of the given $K$; it does not replace it by an outermost
hull.
The volume term is modular under union and intersection, and its
density is zero on $K$. In particular its value on $L\cup K$
equals its value on $L$. Constrained minimality of $K'$ against
$L\cup K$ therefore gives $\mathcal F(K')\le\mathcal F(L)$.
This proves local unconstrained minimality. Its comparison with the
smooth reference metric yields the uniform almost-minimizing estimate
from the H\"older weight, with an additional error at most $Cr^n$
from the bounded volume forcing.
The canonical density representative removes null cracks; for these
almost-minimizing sets, absence of bounded complementary BV components
therefore gives the connected open end exterior asserted above.

\paragraph{Uniform graph estimates before separation.}
At a hypothetical common tangent-cone patch, the local comparison just
proved and one-sided small-excess regularity give uniform
$C^{1,\theta}$ candidate graphs.
For a flow conductor, the original graph and its one-sided metric
have the all-order bounds in
Lemma~\ref{lem:rpi:simultaneous-regularity}. In the smooth initial
setting, the rescaled smooth metric has uniform bounds and the
minimal graph equation supplies the same regular-patch estimates.
Harmonic boundary estimates for $u$, whose original graphical
boundary value is constant, give all fixed finite orders on smaller
patches. The metric $h'$ consequently has controlled extensions
from the original exterior.

For the first graph equation, extend $1-u$ oddly through the original
regular graph. It is negative on the filled side, so the corresponding
$C^{1,\alpha}$ extension of $u$ is at least one there and dominates
its original extension. In the smooth initial setting multiply this
factor into the smooth metric $h$. In the flow setting, first use the
dominating extension of $b^2g_s$ from
Lemma~\ref{lem:rpi:simultaneous-regularity}, and multiply the two
extended conformal factors. After shrinking the patch the resulting
metric is positive, dominates $h'$, and agrees with it on the original
exterior, hence on the candidate graph.
The forcing density is unchanged where $h_0$ is nonzero, because
that region lies in the original exterior; on $K$ it is zero in
both metrics. Thus $K'$ remains a local minimizer of the same
perimeter-plus-volume functional in this regular extension metric.
Its weak Euler equation is exactly
\[
 H_{\nu}^{h'}=-\delta h_0.
\]

The uniformly elliptic prescribed-mean-curvature graph equation
improves the candidate to $C^{2,\alpha}$. The original metric and
harmonic perturbation have the higher-order coefficient extensions
already established, so elliptic bootstrapping gives every fixed
higher-order candidate estimate on smaller patches. After the first
equation is established, these extensions need only match the
one-sided boundary jets; domination is no longer required.
Under a radius-$r_i$ blow-up the forcing is
$r_i(-\delta h_0)(r_i y)$, with uniform local coefficient bounds
of every fixed order. The candidate's estimates therefore precede
and do not assume singular separation.

These estimates already exclude contact at a regular point. At contact with
the original graph, $h_0=0$, whereas the strict harmonic conformal
derivative makes the original graph have strictly negative
end-oriented mean curvature in $h'$. The candidate equation has
zero curvature there. Ordered graph comparison forbids contact.
If only one frontier was initially known to be regular there, the
one-sided regular-neighborhood argument first makes the other
regular. This excludes regular contact only; singular contact is
the separate issue addressed next.

\paragraph{The candidate graph equation and singular detachment.}
Apply Equation~\eqref{eq:rpi:separation-graph-equation} with earlier
metric $g=h$, later metric $h'=u^{4/d}h$, and prescribed curvature
$f=-\delta h_0$. Since $f$ is smooth and zero on $K$, it satisfies
$|f(x)|\le C\operatorname{dist}(x,\Sigma)$. The original distance
from a point at rescaled graph height $\zeta_i$ to $\Sigma$ is at
most $Cr_i\zeta_i$, and therefore
\[
 |U_i^{2/d}f_i|\le Cr_i^2\zeta_i.
\]
Put $b_i=U_i^{2/d}f_i/\zeta_i$ and absorb this term into the operator:
\[
 L_i=J_i-Q_i+b_i,\qquad |b_i|\le Cr_i^2,
 \qquad L_i\zeta_i=F_i\le0.
\]
The coefficient also has a uniform H\"older bound. Indeed, $f$ vanishes
on the earlier graph, so the fundamental theorem of calculus along
the normal graph segment expresses $f_i/\zeta_i$ as $r_i^2$ times
an average of first derivatives of $f$. The smooth forcing and the
uniform graph and metric bounds control that average in the fixed
regular-patch coordinates. Thus $b_i\to0$ and $L_i\to J_C$ with
the coefficient control required by
Proposition~\ref{prop:rpi:singular-separation}.

The sign of $f$ is not needed in this absorption. The strictly
negative right side instead comes from the harmonic conformal factor:
Lemma~\ref{lem:rpi:relative-harmonic}, applied to $1-u>0$, gives the
patchwise two-sided relative bounds for $F_i$. All parameters have
remained fixed during this argument.

\begin{corollary}[Detachment of slices and perturbed enclosures]
\label{cor:rpi:flow-detachment}
For any fixed $0\le s<t$, the conformal-flow frontiers
$\Sigma_s$ and $\Sigma_t$ are disjoint as closed perimeter supports.
For either the flow or smooth initial exterior in
Section~\ref{sec:rpi:perturbed-candidate}, let $u$ satisfy the
harmonic, boundary, strictness, and end hypotheses stated there. For the forcing and
sufficiently small fixed $\delta>0$ constructed there, the
perimeter-plus-volume minimizer is disjoint from the original
frontier, including its singular set. Each fixed pair of these
compact supports has positive distance.
\end{corollary}

\begin{proof}
The flow assertion was proved in
Section~\ref{sec:rpi:flow-slice-separation}. For the perturbed
enclosure, the preceding regular comparison excludes regular
contact. The power almost-minimizing estimate, the regular-patch
metric and graph bounds, and the displayed forcing equation verify
Proposition~\ref{prop:rpi:singular-separation}, which excludes the
remaining singular contact. Compactness of each pair of supports
then gives positive distance. The constants and the resulting
distance may depend on the fixed time pair or perturbation parameter.
\end{proof}
